% sample_chapter.tex -- design-system SAMPLE chapter (not part of any book).
% It shows every environment and macro of hscphys.sty in a realistic setting.
% (sample_chapter_answers.tex is added automatically.)
% hsc-chapter: 2
\chapter{Uniform circular motion}\label{ch2:chap}
\covers{M5.IQ2.01, M5.IQ2.02, M5.IQ2.03}

\inquiry{Why do objects move in circles?}

\begin{goals}
  \item explain why an object moving in a circle at a constant speed is accelerating
  \item use $v = \dfrac{2\pi r}{T}$, $a_\mathrm{c} = \dfrac{v^2}{r}$ and $F_\mathrm{c} = \dfrac{mv^2}{r}$
  \item name the real force that keeps an object in its circle, and test
        $F_\mathrm{c} = \dfrac{mv^2}{r}$ in an experiment
\end{goals}

\begin{syllabus}
  \item M5.IQ2.01: investigate how centripetal force, mass, speed and radius are related
  \item M5.IQ2.02: analyse the forces on objects in uniform circular motion
  \item M5.IQ2.03: solve problems with $a_\mathrm{c}$, $v$, $F_\mathrm{c}$ and $\omega$
\end{syllabus}

%% ------------------------------------------------------------------
\section{Moving in a circle}\label{ch2:sec:circle}

\begin{bigidea}
An object can move at a constant speed and still accelerate. In a circle, the
direction of motion changes all the time. A net force towards the centre causes
this change.
\end{bigidea}

\begin{keywords}
  \kw{Uniform circular motion}{moving in a circle at a constant speed.}
  \kw{Period ($T$)}{the time for one full turn. Unit: second (s).}
  \kw{Centripetal}{pointing towards the centre of the circle.}
  \kw{Centripetal force ($F_\mathrm{c}$)}{the net force towards the centre that
    keeps an object moving in a circle.}
\end{keywords}

Swing a ball on a string in a horizontal circle above your head. The ball goes
round at the same speed all the time. This is called \term{uniform circular
motion}. \emph{Uniform} means ``not changing''. Here the speed does not change.

\begin{refresher}{Velocity is a vector}
Velocity has a size (the speed) and a direction. If the size \emph{or} the
direction changes, the velocity changes. A change in velocity is an
acceleration. Newton's second law, $\vec{F}_\mathrm{net} = m\vec{a}$, says that
an acceleration always needs a net force.
\end{refresher}

\begin{sidebyside}[0.47]
  \centering
  \begin{tikzpicture}[x=1cm, y=1cm]
    \draw[hscmuted!55, line width=0.8pt] (0,0) circle (1.9);
    \draw[->, hscmuted, line width=0.6pt] (140:2.25) arc (140:185:2.25);
    \node[diagramlabel, text=hscmuted, anchor=east] at (162:2.4) {motion};
    \coordinate (A) at (40:1.9);
    \coordinate (B) at (215:1.9);
    \draw[guide] (0,0) -- node[pos=0.3, above left, diagramlabel, inner sep=1pt] {$r$} (A);
    \fill[hscink] (0,0) circle (1.5pt);
    \node[diagramlabel, below right, inner sep=2pt] at (0,0) {centre};
    \draw[velocity] (A) -- ++(130:1.35) node[above, inner sep=2pt, text=hscvel] {$v$};
    \draw[force] (A) -- ($(A)!0.52!(0,0)$) node[pos=0.6, below right, inner sep=1pt, text=hscforce] {$F_\mathrm{c}$};
    \draw[velocity] (B) -- ++(305:1.35) node[right, inner sep=2pt, text=hscvel] {$v$};
    \draw[force] (B) -- ($(B)!0.52!(0,0)$) node[pos=0.55, above left, inner sep=1pt, text=hscforce] {$F_\mathrm{c}$};
    \fill[object] (A) circle (4.5pt);
    \fill[object] (B) circle (4.5pt);
  \end{tikzpicture}
  \captionof{figure}{Top view of a ball on a string at two moments.}
  \label{ch2:fig:topview}
\nextside
  Look at \cref{ch2:fig:topview}.
  \begin{itemize}
    \item The \textcolor{hscvel}{\bfseries velocity} arrow touches the circle.
          It is a \term{tangent} to the circle.
    \item The velocity arrow keeps turning as the ball moves round.
    \item The \textcolor{hscforce}{\bfseries force} arrow points to the centre.
    \item The force is at \ang{90} to the velocity. It changes the direction,
          not the speed.
  \end{itemize}
\end{sidebyside}

What makes the direction change? A net force towards the centre. We call it the
\term{centripetal force} (``centripetal'' means ``centre-seeking'').

The centripetal force is \emph{not} a new kind of force. It is a job that a real
force does. \Cref{ch2:tab:forces} shows which real force does the job.

\begin{table}
  \caption{Which real force gives the centripetal force?}\label{ch2:tab:forces}
  \centering
  \begin{tabular}{@{}L{55mm}L{70mm}@{}}
    \toprule
    \colhead{Situation} & \colhead{Real force towards the centre}\\
    \midrule
    Ball on a string (horizontal circle) & tension in the string\\
    Car on a flat bend & friction between the tyres and the road\\
    Satellite orbiting the Earth & gravity (the Earth's pull)\\
    Car on a banked track & part of the normal force from the road\\
    \bottomrule
  \end{tabular}
\end{table}

If the string breaks, the ball does not fly outwards. With no net force it
moves in a straight line along the tangent (Newton's first law).

\begin{quickcheck}
  \item A runner goes once around a circular track of radius \qty{50}{m} in
        \qty{40}{s}. Calculate her speed.
  \item A car goes around a roundabout at a constant \qty{30}{km.h^{-1}}.
        Explain why it is accelerating.\label{ch2:qc:roundabout}
\end{quickcheck}

%% ------------------------------------------------------------------
\section{Speed, acceleration and force}\label{ch2:sec:formulas}
\covers{M5.IQ2.03}

In one full turn, an object travels the circumference of the circle, $2\pi r$.
It takes one period, $T$, to do this.

\begin{formulacard}{Speed in a circle}
  \cardformula{v = \frac{2\pi r}{T}}
  \begin{symbols}
    \sym{$v$}{speed of the object}{\si{m.s^{-1}}}
    \sym{$r$}{radius of the circle}{\si{m}}
    \sym{$T$}{period: time for one full turn}{\si{s}}
  \end{symbols}
  \inwords{speed = distance for one turn ($2\pi r$) $\div$ time for one turn ($T$).}
  \whentouse{an object goes round a circle at a constant speed.}
  \whennottouse{the speed changes during the turn (for example, a ball whirled
    in a vertical circle).}
  \rearranged{$T = \dfrac{2\pi r}{v}$ \qquad $r = \dfrac{vT}{2\pi}$}
  \tensecond{$r = \qty{2.0}{m}$, $T = \qty{4.0}{s}$, so
    $v = \dfrac{2\pi \times 2.0}{4.0} = \qty{3.1}{m.s^{-1}}$.}
  \trap{use the radius, not the diameter. ``2 turns per second'' means
    $T = \qty{0.50}{s}$ (use $f = \dfrac{1}{T}$).}
\end{formulacard}

\begin{formulacard}{Centripetal force}
  \formula{F_\mathrm{c} = \frac{mv^2}{r}}
  \begin{symbols}
    \sym{$F_\mathrm{c}$}{net force towards the centre}{\si{N}}
    \sym{$m$}{mass of the object}{\si{kg}}
    \sym{$v$}{speed of the object}{\si{m.s^{-1}}}
    \sym{$r$}{radius of the circle}{\si{m}}
  \end{symbols}
  \inwords{a heavier or faster object needs a bigger inward force. A tighter
    circle (smaller $r$) also needs a bigger force.}
  \whentouse{any object in uniform circular motion. It is
    $F_\mathrm{net} = ma$ with $a_\mathrm{c} = \dfrac{v^2}{r}$.}
  \whennottouse{to find a force along the direction of motion.
    $F_\mathrm{c}$ only points to the centre.}
  \rearranged{$v = \sqrt{\dfrac{F_\mathrm{c}\, r}{m}}$ \qquad
    $r = \dfrac{mv^2}{F_\mathrm{c}}$ \qquad $m = \dfrac{F_\mathrm{c}\, r}{v^2}$}
  \tensecond{$m = \qty{2.0}{kg}$, $v = \qty{3.0}{m.s^{-1}}$, $r = \qty{1.5}{m}$,
    so $F_\mathrm{c} = \dfrac{2.0 \times 3.0^2}{1.5} = \qty{12}{N}$.}
  \trap{$F_\mathrm{c}$ is the \emph{net} force. Never draw it as an extra
    arrow next to the tension or the friction.}
\end{formulacard}

\Cref{ch2:fig:forcespeed} shows that the force grows quickly with speed.
Doubling the speed makes the force four times bigger, because $v$ is squared.

\begin{figure}
  \centering
  \includegraphics{sample_force_speed.pdf}
  \caption{Centripetal force against speed for a \qty{2.0}{kg} object, for
    three radii.}\label{ch2:fig:forcespeed}
\end{figure}

Sometimes a question gives the period instead of the speed. You can join the
two formulas into one. Each line below does one small step.
\begin{align}
  a_\mathrm{c} &= \frac{v^2}{r}
     && \text{start with the formula on the sheet}\notag\\
  &= \frac{1}{r}\left(\frac{2\pi r}{T}\right)^{2}
     && \text{replace $v$ with $\dfrac{2\pi r}{T}$}\notag\\
  &= \frac{4\pi^2 r^2}{r\,T^2}
     && \text{square the top and the bottom}\notag\\
  &= \frac{4\pi^2 r}{T^2}
     && \text{cancel one $r$}\label{ch2:eq:acT}
\end{align}
For the ball in \cref{ch2:we:ball} (below), \cref{ch2:eq:acT} gives
$a_\mathrm{c} = \dfrac{4\pi^2 \times 0.80}{0.50^2} = \qty{126}{m.s^{-2}}$,
the same as $\dfrac{v^2}{r}$.

\begin{toolkit}{Radians}
Angles in circular motion are measured in \term{radians} (rad).
One full turn is $2\pi$ rad, which is the same as \ang{360}.
\begin{itemize}
  \item degrees $\to$ radians: multiply by $\dfrac{\pi}{180}$.
        Example: $\ang{90} = 90 \times \dfrac{\pi}{180} = \dfrac{\pi}{2} = \qty{1.57}{rad}$.
  \item Angular velocity $\omega$ is measured in \si{rad.s^{-1}}.
\end{itemize}
\end{toolkit}

\begin{formulacard}[notonsheet]{Speed from angular velocity}
  \cardformula{v = \omega r}
  \begin{symbols}
    \sym{$v$}{speed of a point on the circle}{\si{m.s^{-1}}}
    \sym{$\omega$}{angular velocity: angle turned each second}{\si{rad.s^{-1}}}
    \sym{$r$}{distance from the centre}{\si{m}}
  \end{symbols}
  \inwords{points further from the centre move faster, even though they turn
    through the same angle.}
  \whentouse{you know how fast something spins and need the speed of one point.}
  \whennottouse{$\omega$ is in revolutions per minute or degrees per second.
    Change it to \si{rad.s^{-1}} first.}
  \rearranged{$\omega = \dfrac{v}{r}$ \qquad $r = \dfrac{v}{\omega}$}
  \tensecond{$\omega = \qty{4.0}{rad.s^{-1}}$ and $r = \qty{0.50}{m}$, so
    $v = 4.0 \times 0.50 = \qty{2.0}{m.s^{-1}}$.}
  \derive{start from two formulas on the sheet. In one period the angle is
    $\Delta\theta = 2\pi$.
    \begin{align*}
      \omega &= \frac{\Delta\theta}{t} = \frac{2\pi}{T}
        && \text{one full turn in time $T$}\\
      v &= \frac{2\pi r}{T} = \left(\frac{2\pi}{T}\right) r = \omega r
        && \text{group $\dfrac{2\pi}{T}$ together}
    \end{align*}}
  \trap{$\omega$ must be in \si{rad.s^{-1}}, not in revolutions per minute.}
\end{formulacard}

\begin{workedexample}{A ball on a string}\label{ch2:we:ball}
A \qty{0.15}{kg} ball moves in a horizontal circle of radius \qty{0.80}{m}.
It makes 2.0 full turns every second. Calculate the centripetal force on the ball.
\given{$m = \qty{0.15}{kg}$, \ $r = \qty{0.80}{m}$, \ 2.0 turns each second,
  so $T = \dfrac{1}{2.0} = \qty{0.50}{s}$}
\find{the centripetal force $F_\mathrm{c}$ (first find the speed $v$)}
\formula{$v = \dfrac{2\pi r}{T}$ \quad then \quad $F_\mathrm{c} = \dfrac{mv^2}{r}$}
\substitute{\eqlines{v &= \frac{2\pi \times \qty{0.80}{m}}{\qty{0.50}{s}}
  = \qty{10.05}{m.s^{-1}} \quad\text{(keep the extra digits)}\\
  F_\mathrm{c} &= \frac{\qty{0.15}{kg} \times (\qty{10.05}{m.s^{-1}})^2}{\qty{0.80}{m}}
  = \qty{18.95}{N}}}
\answer{\final{$F_\mathrm{c} = \qty{19}{N}$} towards the centre
  (2 significant figures, like the data).}
\checkstep{units: $\unit{kg} \times (\unit{m.s^{-1}})^2 \div \unit{m}
  = \unit{kg.m.s^{-2}} = \unit{N}$ \cmark. About \qty{19}{N} is a firm pull,
  which fits a ball whirled fast.}
\end{workedexample}

\begin{workedexample}{How fast can a car take a flat bend?}\label{ch2:we:car}
A \qty{1500}{kg} car drives around a flat bend of radius \qty{45}{m}. The largest
friction force the road can give is \qty{8.1e3}{N}. Find the fastest speed the
car can go around the bend without skidding. Give it in \si{m.s^{-1}} and
\si{km.h^{-1}}.
\given{$m = \qty{1500}{kg}$, \ $r = \qty{45}{m}$, \
  largest friction $= \qty{8.1e3}{N}$. Friction gives the centripetal force.}
\find{the largest speed $v_\mathrm{max}$}
\formula{$F_\mathrm{c} = \dfrac{mv^2}{r}$, \ rearranged:
  $v = \sqrt{\dfrac{F_\mathrm{c}\, r}{m}}$}
\substitute{$v_\mathrm{max} = \sqrt{\dfrac{\qty{8.1e3}{N} \times \qty{45}{m}}{\qty{1500}{kg}}}
  = \sqrt{\qty{243}{m^2.s^{-2}}} = \qty{15.6}{m.s^{-1}}$}
\answer{\final{$v_\mathrm{max} = \qty{16}{m.s^{-1}}$}
  \ ($15.6 \times 3.6 = \qty{56}{km.h^{-1}}$)}
\checkstep{\qty{56}{km.h^{-1}} is a sensible speed for a tight town bend.
  Faster than this, the car needs more than \qty{8.1e3}{N}, so it slides to the
  outside of the bend.}
\end{workedexample}

\begin{mistakes}
\begin{itemize}
  \item \xmark\ ``An outward \emph{centrifugal} force pushes the ball out.''
        \cmark\ There is no outward force on the ball. Without the inward force
        it would go straight on, along the tangent.
  \item \xmark\ Using the diameter in $v = \dfrac{2\pi r}{T}$.
        \cmark\ $r$ is the radius.
  \item \xmark\ Putting a speed in \si{km.h^{-1}} into a formula.
        \cmark\ Divide by 3.6 to get \si{m.s^{-1}}.
\end{itemize}
\end{mistakes}

\begin{examtip}
``Explain why an object moving in a circle at constant speed is accelerating''
is a common 3-mark question. Full marks need three linked points:
\begin{enumerate}
  \item the direction of motion keeps changing, so the velocity changes;
  \item a change in velocity is an acceleration, and it points to the centre;
  \item by Newton's second law there is a net force towards the centre. Name
        the real force (tension, friction, gravity, \ldots).
\end{enumerate}
\end{examtip}

\begin{quickcheck}
  \item A \qty{0.50}{kg} ball moves in a circle of radius \qty{1.2}{m} at
        \qty{4.0}{m.s^{-1}}. Calculate the centripetal force.\label{ch2:qc:ball}
  \item The ball in \cref{ch2:qc:ball} now moves twice as fast in the same circle. What
        happens to the centripetal force?
  \item A wheel of radius \qty{0.30}{m} turns through $3\pi$ rad in
        \qty{1.5}{s}. Find its angular velocity and the speed of a point on
        the rim.
\end{quickcheck}

%% ------------------------------------------------------------------
\section{Key experiment: whirling a stopper}\label{ch2:sec:experiment}
\covers{M5.IQ2.01}

\begin{experiment}{Testing $F_\mathrm{c} = mv^2/r$}
\aim{To find how the centripetal force depends on speed (same mass and radius).}
\apparatus{rubber stopper (\qty{0.025}{kg}), string, smooth plastic tube,
  hanging masses (\qtyrange{0.050}{0.250}{kg}), paper clip, metre rule, stopwatch.}
\method{\begin{enumerate}
  \item Thread the string through the tube. Tie the stopper to the top end and
        hang a \qty{0.050}{kg} mass on the bottom end.
  \item Clip a paper clip to the string so the length from the tube to the
        stopper is $r = \qty{0.60}{m}$. Whirl the stopper in a horizontal
        circle, keeping the clip just below the tube.
  \item Time 10 full turns. Divide by 10 to get the period $T$.
  \item Repeat with hanging masses up to \qty{0.250}{kg}.
\end{enumerate}
The hanging mass does not move, so the tension in the string equals its weight
$Mg$. The tension gives the centripetal force: $F_\mathrm{c} = Mg$.}
\results{\begin{center}\small
  \begin{tabular}{@{}cccccc@{}}
    \toprule
    $M$ (\si{kg}) & $F_\mathrm{c} = Mg$ (\si{N}) & time for 10 turns (\si{s})
      & $T$ (\si{s}) & $v$ (\si{m.s^{-1}}) & $v^2$ (\si{m^2.s^{-2}})\\
    \midrule
    0.050 & 0.49 & 11.2 & 1.12 & 3.37 & 11.3\\
    0.100 & 0.98 & 7.7  & 0.77 & 4.90 & 24.0\\
    0.150 & 1.47 & 6.4  & 0.64 & 5.89 & 34.7\\
    0.200 & 1.96 & 5.4  & 0.54 & 6.98 & 48.7\\
    0.250 & 2.45 & 4.9  & 0.49 & 7.69 & 59.2\\
    \bottomrule
  \end{tabular}
  \end{center}
  \begin{figure}
  \centering
  \begin{tikzpicture}
    \begin{axis}[hsc, width=0.62\linewidth, height=0.30\linewidth,
        xmin=0, xmax=65, ymin=0, ymax=2.8,
        xlabel={$v^2$ (\si{m^2.s^{-2}})}, ylabel={$F_\mathrm{c}$ (\si{N})},
        ytick={0,0.5,...,2.5}, yticklabel style={/pgf/number format/.cd, fixed, precision=1}]
      \addplot[hscseriesA, only marks, mark=*, mark size=2pt]
        coordinates {(11.33,0.490) (23.97,0.980) (34.70,1.470) (48.74,1.960) (59.19,2.450)};
      \addplot[hscseriesB, dashed, line width=1pt, domain=0:65, samples=2] {0.04118*x};
      \node[anchor=west, font=\small, text=hscink] at (axis cs:36,0.85)
        {line of best fit: gradient \qty{0.041}{kg.m^{-1}}};
    \end{axis}
  \end{tikzpicture}
  \caption{Centripetal force against speed squared.}\label{ch2:fig:expgraph}
  \end{figure}}
\conclusion{The graph is a straight line through the origin, so
  $F_\mathrm{c} \propto v^2$ when $m$ and $r$ stay the same. The gradient,
  \qty{0.041}{kg.m^{-1}}, is close to
  $\dfrac{m}{r} = \dfrac{0.025}{0.60} = \qty{0.042}{kg.m^{-1}}$.
  This supports $F_\mathrm{c} = \dfrac{mv^2}{r}$.}
\significance{The experiment links a force we can measure (the weight of the
  hanging masses) to the motion (speed and radius). Engineers use the same
  relationship to design road bends and fairground rides.}
\validity{\begin{itemize}
  \item \textbf{Validity:} keep the stopper's mass and the radius the same (use
        the paper clip). Keep the circle as close to horizontal as you can.
  \item \textbf{Reliability:} repeat each trial three times and average the
        times.
  \item \textbf{Accuracy:} time 10 turns, not 1, so your reaction time (about
        \qty{0.2}{s}) matters less. A smooth tube reduces friction.
\end{itemize}}
\end{experiment}

%% ------------------------------------------------------------------
\section{Bringing it together}\label{ch2:sec:summary}

\begin{summarybox}{Uniform circular motion}
\begin{itemize}
  \item Constant speed in a circle is still acceleration: the direction changes.
  \item The net force and the acceleration point to the centre.
  \item $F_\mathrm{c}$ is a job done by a real force (tension, friction, gravity).
  \item Double the speed and you need four times the force.
\end{itemize}
\end{summarybox}

\subsection{Formulas in this chapter}
\begin{formulatable}
  \frow{$v = \dfrac{2\pi r}{T}$}{speed from the radius and the period}
  \frow{$a_\mathrm{c} = \dfrac{v^2}{r}$}{centripetal acceleration (towards the centre)}
  \frow{$F_\mathrm{c} = \dfrac{mv^2}{r}$}{net force needed to keep $m$ in the circle}
  \frow[combine={$a_\mathrm{c} = \frac{v^2}{r}$ and $v = \frac{2\pi r}{T}$}]%
    {$a_\mathrm{c} = \dfrac{4\pi^2 r}{T^2}$}{acceleration when you know the period}
  \frow[notonsheet]{$v = \omega r$}{speed from the angular velocity
    ($\omega$ in \si{rad.s^{-1}})}
\end{formulatable}

\subsection{Words from this chapter}
\begin{glossarytable}
  \gterm{Uniform circular motion}{moving in a circle at a constant speed}{匀速圆周运动}
  \gterm{Centripetal force}{the net force towards the centre of the circle}{向心力}
  \gterm{Period}{the time for one full turn}{周期}
  \gterm{Angular velocity}{the angle turned each second, in \si{rad.s^{-1}}}{角速度}
\end{glossarytable}

%% ------------------------------------------------------------------
\section*{Practice questions (sample)}

\begin{practiceq}{1}{2}
A stone on a string moves in a horizontal circle at a constant speed. Which
statement is correct?
\choices{Its velocity is constant.}{The net force on it is zero.}
        {Its acceleration points to the centre.}{An outward force acts on it.}
\end{practiceq}

\begin{practiceq}{3}{2}
A \qty{1200}{kg} car travels around a flat curve of radius \qty{60}{m} at
\qty{54}{km.h^{-1}}. Calculate the friction force needed.
\end{practiceq}

\begin{practiceq}{3}{2}
Explain why a car on a flat road can skid on an icy bend, even when it is
driven at a speed that is safe on a dry day.
\answerlines{3}
\end{practiceq}
